\documentclass[11pt,a4paper]{article}
\usepackage[utf8]{inputenc}
\usepackage[T1]{fontenc}
\usepackage{lmodern}
\usepackage[margin=1.6cm,top=1.8cm,bottom=1.6cm,headsep=0.5cm]{geometry}
\usepackage{amsmath,amssymb}
\usepackage{enumitem}
\usepackage{titlesec}
\usepackage{fancyhdr}

\pagestyle{fancy}
\fancyhf{}
\lhead{Seconde}
\rhead{Chapitre 2 : Arithmétique dans $\mathbb{N}$ et $\mathbb{Z}$}
\cfoot{\thepage}
\renewcommand{\headrulewidth}{0pt}
\setlength{\parindent}{0pt}
\setlength{\parskip}{0.3em}

\newcommand{\N}{\mathbb{N}}
\newcommand{\Z}{\mathbb{Z}}
\titleformat{\paragraph}{\normalfont\bfseries}{}{0em}{}
\titlespacing*{\paragraph}{0pt}{0.9em}{0.3em}
\newcommand{\exo}[1]{\paragraph{Exercice #1}}
\newcommand{\rep}[1]{\textbf{#1}}
\setlist{itemsep=0.35em, topsep=0.2em, leftmargin=1.8em}

\begin{document}

\begin{center}
  {\Large\bfseries Corrigé de la feuille d'exercices n°2 -- Raisonner en arithmétique}
\end{center}

\exo{1}
Pour montrer qu'une affirmation est fausse, il suffit d'un contre-exemple.
\begin{enumerate}[label=\alph*)]
  \item $2$ et $2$ sont pairs, mais $2 \times 2 = 4$ n'est pas un multiple de $8$.
  \item $12 = 4 \times 3$ et $12 = 6 \times 2$ : $12$ est multiple de $4$ et de $6$. Mais $12$ n'est pas un multiple de $24$ (car $0 < 12 < 24$).
  \item Pour $n = 4$ : $4^2 + 4 + 1 = 21 = 3 \times 7$, qui n'est pas un nombre premier.
  \item Pour $a = 2$, $b = 1$ et $c = 1$ : $b + c = 2$ et $2$ divise $2$. Mais $2$ ne divise pas $1$ : $a$ ne divise ni $b$, ni $c$.
\end{enumerate}

\exo{2}
\rep{Affirmation A : vraie.} Soient $n$ et $n+1$ deux entiers consécutifs, avec $n \in \Z$. Alors
\[ n + (n+1) = 2n + 1. \]
Or $n \in \Z$, donc $n + (n+1)$ est impair.

\rep{Affirmation B : fausse.} Contre-exemple : $2 + 3 + 4 = 9$, qui n'est pas un multiple de $6$.

\exo{3}
\begin{enumerate}[label=\rep{\arabic*.}]
  \item On étudie les deux cas possibles.
  \begin{itemize}
    \item Si $n$ est pair : il existe $k \in \Z$ tel que $n = 2k$. Alors
    \[ n^2 - n = 4k^2 - 2k = 2(2k^2 - k). \]
    Or $2k^2 - k \in \Z$, donc $n^2 - n$ est pair.
    \item Si $n$ est impair : il existe $k \in \Z$ tel que $n = 2k + 1$. Alors
    \[ n^2 - n = (2k+1)^2 - (2k+1) = 4k^2 + 4k + 1 - 2k - 1 = 4k^2 + 2k = 2(2k^2 + k). \]
    Or $2k^2 + k \in \Z$, donc $n^2 - n$ est pair.
  \end{itemize}
  Dans tous les cas, $n^2 - n$ est pair.

  \item On factorise : $n^2 + 5n = n(n+5)$.
  \begin{itemize}
    \item Si $n$ est pair : $n = 2k$ avec $k \in \Z$, donc $n(n+5) = 2k(2k+5) = 2 \times k(2k+5)$.
    \item Si $n$ est impair : $n = 2k + 1$ avec $k \in \Z$, donc $n + 5 = 2k + 6 = 2(k+3)$ et
    $n(n+5) = 2 \times (2k+1)(k+3)$.
  \end{itemize}
  Dans tous les cas, $n^2 + 5n$ est pair : il existe $m \in \Z$ tel que $n^2 + 5n = 2m$. Alors
  \[ n^2 + 5n + 7 = 2m + 6 + 1 = 2(m+3) + 1. \]
  Or $m + 3 \in \Z$, donc $n^2 + 5n + 7$ est impair.
\end{enumerate}

\exo{4}
Soit $k \in \Z$.
\begin{enumerate}[label=\rep{\arabic*.}]
  \item On étudie les trois cas possibles.
  \begin{itemize}
    \item Si $n = 3k$ : $n(n+1)(n+2) = 3k(3k+1)(3k+2) = 3 \times k(3k+1)(3k+2)$.
    \item Si $n = 3k + 1$ : $n + 2 = 3k + 3 = 3(k+1)$, donc $n(n+1)(n+2) = 3 \times (3k+1)(3k+2)(k+1)$.
    \item Si $n = 3k + 2$ : $n + 1 = 3k + 3 = 3(k+1)$, donc $n(n+1)(n+2) = 3 \times (3k+2)(k+1)(3k+4)$.
  \end{itemize}
  Dans chaque cas, le facteur qui multiplie $3$ est un produit d'entiers, donc un entier.
  Dans tous les cas, $n(n+1)(n+2)$ est un multiple de $3$.

  \item On calcule $n^2 + 1$ dans les trois cas.
  \begin{itemize}
    \item Si $n = 3k$ : $n^2 + 1 = 9k^2 + 1 = 3 \times 3k^2 + 1$.
    \item Si $n = 3k + 1$ : $n^2 + 1 = 9k^2 + 6k + 1 + 1 = 3(3k^2 + 2k) + 2$.
    \item Si $n = 3k + 2$ : $n^2 + 1 = 9k^2 + 12k + 4 + 1 = 3(3k^2 + 4k + 1) + 2$.
  \end{itemize}
  Dans tous les cas, $n^2 + 1$ s'écrit $3m + 1$ ou $3m + 2$ avec $m \in \Z$ : il est strictement compris entre les deux multiples de $3$ consécutifs $3m$ et $3m + 3$.
  Donc $n^2 + 1$ n'est jamais un multiple de $3$.
\end{enumerate}

\exo{5}
\begin{enumerate}[label=\rep{\arabic*.}]
  \item On suppose que $n^2$ est impair et que $n$ est pair. Il existe alors $k \in \Z$ tel que $n = 2k$, et
  \[ n^2 = 4k^2 = 2 \times 2k^2. \]
  Or $2k^2 \in \Z$, donc $n^2$ est pair : c'est une contradiction. Donc, si $n^2$ est impair, $n$ est impair.

  \item On suppose qu'il existe des entiers $a$ et $b$ tels que $2a + 4b = 7$. Alors
  \[ 7 = 2(a + 2b). \]
  Or $a + 2b \in \Z$, donc $7$ serait pair. C'est une contradiction, car $7 = 2 \times 3 + 1$ est impair.
  Il n'existe donc pas d'entiers $a$ et $b$ tels que $2a + 4b = 7$.

  \item On suppose qu'il existe des entiers $a$ et $b$ tels que $6a + 9b = 10$. Alors
  \[ 10 = 3(2a + 3b). \]
  Or $2a + 3b \in \Z$, donc $10$ serait un multiple de $3$. C'est une contradiction, car $\dfrac{10}{3}$ n'est pas un entier ($10 = 3 \times 3 + 1$).
  Il n'existe donc pas d'entiers $a$ et $b$ tels que $6a + 9b = 10$.
\end{enumerate}

\exo{6}
\begin{enumerate}[label=\rep{\arabic*.}]
  \item $n^2 - 1 = (2k+1)^2 - 1 = 4k^2 + 4k + 1 - 1 = 4k^2 + 4k = 4k(k+1)$.

  \item D'après la question 1 de l'exercice 3, $m^2 - m$ est pair pour tout $m \in \Z$. Avec $m = k + 1$ :
  \[ (k+1)^2 - (k+1) = (k+1)\big((k+1) - 1\big) = (k+1)k. \]
  Donc $k(k+1)$ est pair.

  \item $k(k+1)$ est pair : il existe $p \in \Z$ tel que $k(k+1) = 2p$. Alors
  \[ n^2 - 1 = 4k(k+1) = 4 \times 2p = 8p. \]
  Or $p \in \Z$, donc $n^2 - 1$ est un multiple de $8$.

  Vérification : $3^2 - 1 = 8 = 8 \times 1$ \quad ; \quad $5^2 - 1 = 24 = 8 \times 3$ \quad ; \quad $7^2 - 1 = 48 = 8 \times 6$.
\end{enumerate}

\end{document}
